\section{Quo Vadis and the Inverted Cross}

The narrative everyone knows---Peter fleeing Rome, meeting Christ on the
road, turning back to be crucified upside down---comes from a single
identifiable source: the apocryphal \work{Acts of Peter}, a Greek romance
that M.~R. James, whose 1924 translation is used here, judged ``written,
probably by a resident in Asia Minor (he does not know much about Rome),
not later than A.D.\ 200.'' It is a work of devotional fiction: most of its
bulk is a wonder-contest at Rome between Peter and Simon Magus (grown from
the Acts 8 episode), ending when Simon, flying over the city, is brought
down by Peter's prayer. Its martyrdom chapters, however, preserved the two
images that fixed Peter's death in Christian imagination, and they deserve
exact quotation with their real character labeled.

The \textit{Quo vadis} scene: persuaded by the brethren to escape a
persecution provoked---in this telling---not by Nero's fire but by
converted concubines of the prefect Agrippa, Peter is leaving the city when
``he saw the Lord entering into Rome. And when he saw him, he said: Lord,
whither goest thou thus? And the Lord said unto him: I go into Rome to be
crucified. And Peter said unto him: Lord, art thou being crucified again?
He said unto him: Yea, Peter, I am being crucified again. And Peter came to
himself\ldots\ and he returned to Rome, rejoicing, and glorifying the Lord''
(\work{Acts of Peter} 35). Condemned by Agrippa ``on an accusation of
godlessness,'' Peter at the cross asks the executioners: ``I beseech you
the executioners, crucify me thus, with the head downward and not
otherwise'' (37). The romance's own motive is not humility but mystical
symbolism: hanging inverted, Peter expounds ``the mystery of all nature,''
for ``the first man, whose race I bear in mine appearance, fell head
downwards'' and so inverted the order of creation (38). The familiar
humility motive is a later smoothing, first extant in patristic summary:
Jerome's notice that Peter was crucified head-down ``asserting that he was
unworthy to be crucified in the same manner as his Lord'' (\work{On
Illustrious Men} 1). None of this is first-century evidence; all of it is
evidence of how the second and fourth centuries venerated a death whose
fact they had received.

Around the same figure other traditions accreted, each with its own
ceiling. Jerome's same notice transmits the developed chronology: Peter
bishop of Antioch, then to Rome ``in the second year of Claudius'' (42) ``to
overthrow Simon Magus,'' holding ``the sacerdotal chair there for
twenty-five years.'' The twenty-five-year Roman episcopate is a
third-century construction---the 1912 Catholic Encyclopedia, no hostile
witness, concedes ``we are unable to trace its origin''---and the Claudian
arrival date rides on the Simon Magus legend; neither is usable chronology,
and both collide with Peter's documented presence in Jerusalem through the
council of c.~48--49. The Mamertine imprisonment with its miraculous
spring, the martyr Petronilla reimagined as Peter's daughter, and the
relic-chair of the \textit{Cathedra Petri} enshrined by Bernini are later
Roman traditions of cult, not biography. The liturgical date itself
belongs to reception history: 29 June is attested from the year 258 as the
day of a commemoration of both apostles connected with the catacombs
\textit{ad Catacumbas}, and became the immovable joint feast; no ancient
source establishes it as the calendar day on which Peter died. Each of
these items is tabulated with its earliest witness in the tradition audit
appendix; none is load-bearing in the historical chapters of this study.
